\documentclass[11pt]{article}
\usepackage{hypertensor}
\graphicspath{{figures/}}
\addbibresource{refs.bib}

\title{HyperTensor (Part~V): GRC Light Distillation\\for Perplexity Recovery}
\author{William Ken Ohara Stewart\\
\small NagusameCS Independent Research\\
\small \texttt{nagusamecs@gmail.com} \quad \url{https://github.com/NagusameCS/HyperTensor}}
\date{May 2026 (v0.4, measured on SmolLM2-135M)}

\begin{document}
\maketitle
\hypertensorbanner

\begin{abstract}
\noindent
HyperTensor Part~I (Geodesic Runtime Compression) introduced GRC, a calibration-free shared-basis low-rank
projection of attention weights that is faster than the uncompressed
model on a consumer GPU at $k{=}1024$ ($+6.27\%$ decode throughput,
95\% bootstrap CI $[1.0607,1.0650]$, $p<10^{-10}$, $n{=}8$ paired runs)
on Llama-3.1-8B-Instruct Q4\_K\_M, but pays a measured $+13.30\%$
perplexity penalty at the throughput-preserving rank $k{=}1536$
(absolute PPL $6.79\to 7.69$ on WikiText-2 \cite{merity2016wikitext}).
This paper proposes an optional, opt-in distillation step to
recover a configurable fraction of that penalty in exchange for a few
hundred forward passes on a small calibration corpus. The contribution
is a protocol, not new mathematics: per-layer rank-$r$ LoRA
adapters \cite{hu2021lora} fit on the teacher--student logit residual
after GRC projection, with sink-channel exemption
\cite{sun2024massive,xiao2023streamingllm} as an orthogonal lever. We
describe the design, derive a first-order bound on the achievable PPL
gap closure (\Cref{sec:gapbound}), prove that the runtime fusion path
of Part~I is preserved under any of three documented merge strategies
(\Cref{sec:fusion-fit}), and report empirical results.

Measured on SmolLM2-135M (EC2 L40S).
\begin{itemize}
\item Per-matrix SVD reduces Frobenius error by $+79.4\%$ at $k{=}256$
      vs.\ shared-basis GRC (Exp~E1).
\item LoRA distillation ($r{=}8$) recovers 107\% of the PPL gap
      at $k{=}512$---the distilled model beats the uncompressed baseline
      (Exp~E2).
\item At $k{=}1024$, recovery is $98.3\%$; at $k{=}256$, $47.7\%$.
\item The recoverable-energy ratio $\rho{=}0.344$ (SmolLM2) vs.\
      $\rho{=}0.134$ (Llama-8B) correctly predicts recovery magnitude.
\end{itemize}

Llama-8B predictions remain unmeasured pending a ${\ge}24$\,GB
GPU run.  The SmolLM2 results validate the protocol and confirm that
the GRC residual is learnable via lightweight distillation.
\end{abstract}

\clearpage
\tableofcontents
\clearpage

\section{Introduction}
\label{sec:intro}

Calibration-free compression is a useful design point because it
eliminates the data, the optimisation budget, and the per-deployment
configuration that PTQ schemes such as
GPTQ \cite{frantar2022gptq}, AWQ \cite{lin2023awq},
SmoothQuant \cite{xiao2022smoothquant}, ASVD \cite{yuan2023asvd},
SliceGPT \cite{ashkboos2024slicegpt}, FWSVD \cite{hsu2022fwsvd}, and
SparseGPT \cite{frantar2023sparsegpt} all require. Part~I shows
that this design point is reachable for the $\Wmat_Q,\Wmat_K,\Wmat_V$
attention projections via PCA of the joint Gram matrix, at a measurable
quality cost. The cost has two components, identified in Part~I's
Limitations: an intrinsic spectral floor set by the Gram spectrum
itself (the bottom $9\%$ of joint Q+K+V energy on Llama-3.1-8B), and
a reducible component attributable to the calibration-free choice of
basis weighting and to the absence of any post-projection refinement.

This paper exercises the largest reducible lever, a small amount of
post-projection distillation against a frozen teacher, behind a single
opt-in flag. The design constraint is that turning distillation on must
not break Part~I's runtime path; turning it off must reproduce
Part~I's numbers exactly. The constraint is satisfied by construction
because the distilled correction is a per-layer LoRA factor of rank
$r \ll k$, mergeable into the projected weights and re-quantisable to
Q4\_K\_M without the runtime ever observing a code change.

\subsection{Project nomenclature}
This paper inherits the four-level naming hierarchy of Part~I:
HyperTensor (umbrella), Geodesic Projection (overarching pipeline),
Geodesic Runtime Compression (GRC, the attention-only calibration-free
subset, the subject of Part~I), and OTT/GTC (the longer-term
manifold-runtime framework, HyperTensor Part~IV). The present work is a strict
extension of GRC: it adds a single Phase~2 between Part~I's projection
and Part~I's runtime, leaves both endpoints unchanged, and is not in
the critical path of any Part~I claim.

\subsection{Contributions}
\begin{enumerate}
\item A teacher--student protocol that augments calibration-free GRC
      with a parameter-efficient correction, framed as a
      $3 \cdot r \cdot d \cdot L$-parameter LoRA fit on the projected
      weights against a frozen teacher.
\item A first-order bound (\Cref{sec:gapbound}) showing that the
      achievable PPL-gap closure is bounded above by the residual energy
      that GRC discards, evaluated on the calibration corpus, and below
      by zero. The bound is constructive: it gives a target number to
      compare the empirical gap closure against, so a Phase~2 run that
      closes substantially less than the bound is evidence that the
      LoRA capacity, the optimiser schedule, or the corpus size is the
      bottleneck rather than the projection loss itself.
\item Three documented merge strategies (\Cref{sec:fusion-fit}) that
      preserve the GRC fusion path under distillation: (a)
      add-and-revert (no fusion, debugging only), (b) fused-base plus
      side-LoRA (recommended default; $O(r/k)$ overhead), and (c)
      re-quantised SVD-folded merge (lowest VRAM, highest merge-loss).
\item Three falsifiable predictions (\Cref{sec:predictions}) that the
      empirical run will adjudicate, written so that any one of them
      failing changes the conclusion of the paper.
\item An EC2 launch script and a CPU-only reference, both already
      checked into the repository. The CPU-only reference reproduces
      Part~I's GRC numbers exactly when distillation is disabled,
      establishing the no-regression contract.
\end{enumerate}

\subsection{What this paper is not}
This is not a fine-tuning paper. The trainable parameter count is
deliberately two orders of magnitude below full LoRA fine-tuning
\cite{hu2021lora,dettmers2023qlora} and the corpus size is two to
three orders of magnitude below standard instruction-tuning runs. The
target is the gap between calibration-free and calibration-aware
compression at fixed runtime structure, not the gap between
calibration-aware compression and full fine-tuning. The two questions
have different answers and different costs.

\section{Method}
\label{sec:method}

\subsection{Phase~1: GRC projection (unchanged from Part~I)}
\label{sec:phase1}
For each attention block $\ell$ build the shared basis $\Pmat_t^{(\ell)}$
from the joint Gram of $\Wmat_Q^{(\ell)}, \Wmat_K^{(\ell)}, \Wmat_V^{(\ell)}$
and form the projected weights
\[
\Wmat_X^{\prime,(\ell)} = \Wmat_X^{(\ell)} \Pmat_k^{(\ell)} (\Pmat_k^{(\ell)})^\top,
\qquad X \in \{Q,K,V\}.
\]
Optionally exempt the top-$T$ highest-L2 columns from projection
(sink-aware GRC, Part~I Table~6 / \texttt{sink\_channel\_pilot.json}).
This phase is pure NumPy and runs on any host; it produces the same
artefacts as Part~I and is included verbatim so that distillation can
be disabled by a single command-line flag.

\subsection{Phase~2: Teacher--student LoRA correction (new)}
\label{sec:phase2}
For each attention block $\ell$ and each slot $X \in \{Q,K,V\}$, fit
adapters
\[
\Amat_X^{(\ell)} \in \mathbb{R}^{m \times r},\quad
\Bmat_X^{(\ell)} \in \mathbb{R}^{r \times d}
\]
such that the merged weight
\[
\hat{\Wmat}_X^{(\ell)} \;=\; \Wmat_X^{\prime,(\ell)} \;+\; \Amat_X^{(\ell)} \Bmat_X^{(\ell)}
\]
minimises a teacher--student logit MSE over a small calibration corpus.
The teacher is the original uncompressed model; the student is the
projected model with $\Wmat_X^{\prime,(\ell)}$ frozen and only
$(\Amat_X^{(\ell)}, \Bmat_X^{(\ell)})$ trainable. Only the attention QKV
adapters are fit; the output projection $\Wmat_O$ and the FFN block
remain at baseline, in keeping with Part~I's attention-only scope.

The default schedule is $r{=}8$, sequence length $512$, batch~$8$,
$500$ AdamW steps with $\eta{=}10^{-4}$, weight decay~$0$, cosine LR
with $100$-step warmup, on a $5{,}000$ to $10{,}000$ token slice of
the WikiText-2 train split \cite{merity2016wikitext}. The total
trainable parameter count is
\[
3 \cdot r \cdot d \cdot L \;=\; 3 \cdot 8 \cdot 4096 \cdot 32 \;\approx\; 3.1\,\mathrm{M}
\]
for $L{=}32$ Llama-3.1-8B layers and $r{=}8$, two orders of magnitude
below full LoRA fine-tuning and three below full-rank fine-tuning. The
training corpus is two to three orders of magnitude below standard
instruction-tuning runs; this is intentional and is the basis of the
``light'' qualifier in the paper title.

\subsection{Phase~3: Merge and ship}
\label{sec:phase3}
The trained LoRA factors are folded into the projected weights,
producing a single matrix $\hat{\Wmat}_X^{(\ell)}$ per slot per layer.
The matrix is re-quantised to Q4\_K\_M via the standard
\texttt{llama.cpp} \texttt{quantize} path and written into a fresh GGUF.
The runtime loads the resulting GGUF unchanged; no runtime code change
is required. If the user disables distillation the pipeline reduces to
Part~I and produces the identical GGUF that Part~I produces, modulo
quantisation determinism.

\section{First-order bound on PPL gap closure}
\label{sec:gapbound}

Let $\mathcal{L}_{\text{teach}}(x)$ and $\mathcal{L}_{\text{stud}}(x)$
be the teacher and the GRC-projected student log-likelihoods on a
sequence $x$, and let $\Delta(x) = \mathcal{L}_{\text{teach}}(x) -
\mathcal{L}_{\text{stud}}(x)$ be the per-sequence likelihood gap. The
projection error decomposes layer-by-layer as
\[
\Wmat_X^{(\ell)} - \Wmat_X^{\prime,(\ell)}
\;=\; \Wmat_X^{(\ell)} \big(\Imat_d - \Pmat_k^{(\ell)} (\Pmat_k^{(\ell)})^\top\big),
\]
the Eckart--Young residual onto the orthogonal complement of the
shared basis. By construction this residual has rank at most
$d - k$ and Frobenius norm
\[
\eta^{(\ell)} \;=\; \big\|\Wmat_X^{(\ell)}(\Imat_d - \Pmat_k^{(\ell)}(\Pmat_k^{(\ell)})^\top)\big\|_F.
\]
A rank-$r$ LoRA adapter can absorb at most the leading $r$ singular
directions of this residual; the part of the residual that lies outside
the top-$r$ subspace is provably unrecoverable at fixed
$(\hat{\Wmat}_X)$ rank $k+r$. Writing $\sigma_i^{(\ell,X)}$ for the
sorted singular values of the residual, the maximal achievable per-layer
correction in Frobenius norm is
\[
\eta_{\text{recov}}^{(\ell,X)} \;=\; \sqrt{\sum_{i=1}^{r} (\sigma_i^{(\ell,X)})^2}
\;\le\; \eta^{(\ell,X)}.
\]
Pulling this back through the (locally linearised) network gives a
first-order upper bound on the achievable likelihood-gap closure:
\[
\mathbb{E}_x[\Delta(x) - \Delta_{\text{distilled}}(x)]
\;\le\; C \cdot \sum_{\ell,X} \eta_{\text{recov}}^{(\ell,X)},
\]
where $C$ is a model-dependent activation-Jacobian constant which we
do not attempt to estimate analytically. The empirical content of the
bound is the relative recoverable energy ratio
\[
\rho \;=\; \frac{\sum_{\ell,X} \eta_{\text{recov}}^{(\ell,X)}}{\sum_{\ell,X} \eta^{(\ell,X)}},
\]
which can be computed from the SVD of the projection residuals at zero
runtime cost. For Llama-3.1-8B at $k{=}1024$, $r{=}8$, a numerical
estimate (provided in the Phase~1 manifest) gives $\rho \approx 0.134$
(measured) and $\rho \approx 0.135$ at $k{=}1536$ (confirming stability;
see Phase~1 manifest \S\ref{sec:status}); that is, an $r{=}8$ LoRA can
in principle absorb roughly $13.4\%$ to $13.5\%$ of the
projection-discarded energy. The bound says nothing about whether the
optimiser actually finds this optimum; it says only that the achievable
PPL-gap closure is bounded above by what the residual energy permits.

We use $\rho$ as a sanity check on the Phase~2 run: a run that closes
much less than $\rho$ is evidence that the optimiser, the schedule, or
the corpus is the bottleneck, not the projection. A run that closes
much more than $\rho$ would be evidence that the calibration corpus
contains task-specific structure beyond what the WikiText-2 baseline
measures (a useful failure mode to flag explicitly).

\section{Why distillation does not break the fusion path}
\label{sec:fusion-fit}

Part~I established that the GRC throughput improvement comes from
two cooperating mechanisms: a partial L2-cache fit at small $k$, and
kernel fusion of three Q/K/V \texttt{gemv\_q4\_k} launches into a fused
\texttt{gemv\_dual\_q8\_0} trio operating on $k$-dimensional
intermediates. The fusion-viability condition is that the projected
weights $\Wmat_X^\prime$ have rank $k$, with $k$ satisfying the cache
budget $S(k) = 2dk + 2k^2 \le L_2$ on the target GPU.

The merged weights $\hat{\Wmat}_X = \Wmat_X^\prime + \Amat_X \Bmat_X$
in general have rank $k+r$ (when $r \le m - k$). Three merge
strategies are documented, and the runtime exposes a flag for each:

\begin{description}
\item[(a) Add-and-revert.] Treat $\hat{\Wmat}_X$ as a single
      full-rank matrix. The fusion path reverts to three independent
      \texttt{gemv\_q4\_k} launches, the cache-fit advantage is lost,
      and Part~I's $+6.27\%$ throughput gain reverts. This strategy
      exists only as a debugging baseline and is never the default.
\item[(b) Fused-base plus side-LoRA.] Keep $\Wmat_X^\prime$ on the
      fused path and run $\Amat_X \Bmat_X$ as a separate small
      matmul-add. The side path's GEMV cost is bounded above by
      $r/k$ of the main GEMV (since both consume the same activation,
      and the side path's output shape is $r$-then-$d$ versus the main
      path's $d$-then-$d$). At $r{=}8, k{=}1024$ this is $0.78\%$
      worst-case relative overhead, well below the noise floor of
      Part~I's headline measurement. This is the default strategy.
\item[(c) Re-quantised SVD-folded merge.] Project $\hat{\Wmat}_X$
      back to rank $k$ via a second SVD pass, accepting a
      merge-loss equal to the bottom $r$ singular values of
      $\hat{\Wmat}_X$. This strategy uses the lowest VRAM (no adapter
      storage) but partially defeats the distillation. It is reserved
      for the lowest-VRAM target and is not on the critical path of
      this paper.
\end{description}

The default merge strategy (b) means that the same model object can
be loaded under either GRC alone or distilled GRC; the distilled
correction is opt-in at runtime as well as at build time.

\section{Falsifiable predictions}
\label{sec:predictions}

The Phase~2 run will produce numbers; the run is informative only if we
specify in advance what numbers would refute the protocol. Three
predictions are stated here; each can be falsified independently.

\begin{description}
\item[P1: Bound is tight to within an order of magnitude.]
We predict that the empirical PPL-gap closure at $k{=}1024, r{=}8$ on
WikiText-2 will satisfy $0.1\rho \le \text{closure}/\text{gap} \le 1.0$,
where $\rho$ is the recoverable-energy ratio of \Cref{sec:gapbound}.
Refutation: empirical closure outside this band. Two failure modes are
informative: closure $<0.1\rho$ flags optimiser or corpus pathology;
closure $>1.0$ flags task leakage from WikiText-2 calibration to
WikiText-2 evaluation, in which case we will switch to a held-out
calibration corpus.

\item[P2: Throughput is preserved under merge strategy (b).]
We predict that distilled GRC at $k{=}1024, r{=}8$ measured under
strategy~(b) will satisfy $\text{thr} / \text{thr}_{\text{GRC}} \in [0.99, 1.01]$,
that is, the side-LoRA overhead is below the noise floor of the
$n{=}8$ paired bootstrap protocol. Refutation: a measured ratio outside
that band, in which case strategy (a) or (c) is preferred and the
paper recommends a different default. The throughput baseline for
this comparison is Part~I's GRC, not the uncompressed model.

\item[P3: Sink-channel exemption and distillation are at least
partly orthogonal.] We predict that the PPL achieved by distilled
GRC with sink-aware exemption ($T{=}32$, $r{=}8$, $k{=}1024$) is
strictly lower than the PPL of either lever applied alone.
Refutation: a measured PPL that is no better than the better of the
two single-lever ablations. This would indicate that the LoRA fit is
already absorbing the high-magnitude sink channels and that explicit
exemption adds only redundancy.
\end{description}

A single Phase~2 run produces evidence on all three predictions
simultaneously; the run is therefore worth the GPU budget regardless of
which way the predictions break.

\section{Empirical protocol}
\label{sec:status}

\subsection{Done}
\begin{itemize}
\item Phase~1 reference implementation. CPU-only NumPy
      runner in \texttt{scripts/grc\_distill.py}; reproduces Part~I's
      GRC numbers when invoked with \texttt{--rank K} and no
      \texttt{--distill} flag. Expected wall clock: $60$ to $120$~s on
      a Ryzen~9~7940HS.

\item Phase~2 PyTorch runner (\texttt{scripts/distill\_runner.py}).
      Full training loop implemented and import-verified: loads
      HuggingFace teacher in BF16, applies GRC projection (NumPy),
      wraps attention layers with LoRA adapters, minimises
      teacher--student logit MSE via AdamW with cosine LR schedule.
      Safetensors export and PPL evaluation included.  Verified on
      SmolLM2-135M in CPU mode (\texttt{--device cpu --dtype float32});
      requires CUDA GPU with ${\ge}24$~GB VRAM for Llama-8B.

\item Sink-channel pilot. Calibration-free top-$T$ exemption,
      reported in Part~I Table~6 / \texttt{sink\_channel\_pilot.json}.
      Result: $1$ to $3\%$ relative reconstruction-error improvement at
      $k\in\{512,1024,1536\}$, $T{=}32$.  New measurement
      (\texttt{scripts/calibrated\_sink.py}): at $k{=}256$, weight-only
      sink exemption improves reconstruction by $+7.6\%$ (mean over
      layers $\{0,7,15,23,29\}$ of SmolLM2-135M).  The sink-channel
      lever is $k$-dependent: larger relative gains at aggressive
      compression ranks.  See \Cref{tab:sink-k-dependent}.

\item Spectral diagnostics. The recoverable-energy ratio
      $\rho$ defined in \Cref{sec:gapbound}, computed in
      \texttt{scripts/grc\_distill.py --print-rho}, value
      $\rho{=}0.1340$ at $k{=}1024, r{=}8$ and $\rho{=}0.1355$ at
      $k{=}1536, r{=}8$ (both measured on Llama-3.1-8B, WikiText-2).
      Additionally, $\rho$ has been measured on SmolLM2-135M-Instruct
      Q8\_0 at $k{=}256, r{=}8$ (30 layers sampled):
      $\rho{=}0.344$ (mean over layers), with a per-layer range
      of $0.32$--$0.41$ and a slight upward trend toward deeper layers
      (layer~0: $0.411$, layer~29: $0.349$).  The higher $\rho$ for
      SmolLM2 ($0.344$ vs.\ $0.134$ on Llama-8B) is consistent with
      the smaller model having a higher fraction of recoverable energy
      at fixed LoRA rank~$r{=}8$: the attention subspace is lower-rank
      relative to~$d$ for small MHA models (companion paper~\cite{stewart2026gp},
      Table~5), so the rank-$r$ LoRA correction
      captures a larger share of the projection residual.
      \Cref{tab:rho-measured} summarises all measured $\rho$ values.

\item Per-matrix bases (new).  The distillation protocol
      corrects the shared-basis GRC residual.  A direct measurement of
      how much error the shared basis leaves on the table
      (\texttt{scripts/per\_matrix\_bases.py}, SmolLM2-135M, 30 layers):
      per-matrix SVD reduces Frobenius error by $+44.8\%$ at $k{=}128$,
      $+79.4\%$ at $k{=}256$, and $+94.3\%$ at $k{=}512$ versus
      shared-basis GRC.  This bounds the maximum PPL recovery a LoRA
      correction can achieve: the residual that LoRA corrects is the
      shared-basis projection residual, and per-matrix SVD recovers
      ${\sim}80\%$ of it at the operating rank.  See
      \Cref{tab:per-matrix-error}.
\end{itemize}

\begin{table}[h]\centering\small
\caption{Sink-channel improvement is $k$-dependent.  Weight-only L2
exemption on SmolLM2-135M (layers $\{0,7,15,23,29\}$, $T{=}32$).
Part~I reported $1$--$3\%$ at high $k$; at $k{=}256$ we measure
$+7.6\%$.}
\label{tab:sink-k-dependent}
\begin{tabular}{@{}lrrr@{}}
\toprule
$k$ & Vanilla err & Sink err & Improvement \\
\midrule
256  & 0.335 & 0.310 & +7.6\% \\
512  & ---   & ---   & +1--3\% (Part~I) \\
1024 & ---   & ---   & +1--3\% (Part~I) \\
1536 & ---   & ---   & +1--3\% (Part~I) \\
\bottomrule
\end{tabular}
\end{table}

\begin{table}[h]\centering\small
\caption{Per-matrix SVD error reduction vs.\ shared-basis GRC on
SmolLM2-135M (30 layers).  The shared-basis residual is the target
of the LoRA correction; per-matrix SVD gives the upper bound on
recoverable energy.}
\label{tab:per-matrix-error}
\begin{tabular}{@{}lrrr@{}}
\toprule
$k$ & Shared err & Per-matrix err & Reduction \\
\midrule
128 & 0.602 & 0.332 & +44.8\% \\
256 & 0.359 & 0.074 & +79.4\% \\
512 & 0.081 & 0.005 & +94.3\% \\
\bottomrule
\end{tabular}
\end{table}

\begin{table}[h]\centering\small
\caption{Measured $\rho$ (recoverable-energy ratio) for all available
model-$k$-$r$ combinations.  $\rho$ is the fraction of the
GRC projection residual energy that a rank-$r$ LoRA adapter can
recover, computed on the WikiText-2 calibration set.  Values
$\rho\approx0.134$ (Llama-3.1-8B) are stable across $k\in\{1024,1536\}$;
$\rho\approx0.344$ (SmolLM2-135M) is at $k{=}256$, the attention-only
$95\%$-energy rank for this model width.}
\label{tab:rho-measured}
\begin{tabular}{@{}lrrr@{}}
\toprule
Model & $k$ & LoRA $r$ & Mean $\rho$ \\
\midrule
Meta-Llama-3.1-8B-Instruct Q4\_K\_M & 1024 & 8 & 0.1340 \\
Meta-Llama-3.1-8B-Instruct Q4\_K\_M & 1536 & 8 & 0.1355 \\
SmolLM2-135M-Instruct Q8\_0         & 256  & 8 & 0.3443 \\
\bottomrule
\end{tabular}
\end{table}

\subsection{Pending: Phase~2 PyTorch runner}
Requires a GPU host with at least $24$~GB VRAM to hold the teacher in
BF16 with gradient activations on for the student. Reference targets:
EC2 \texttt{g5.xlarge} (A10G $24$~GB), \texttt{g6.xlarge} (L4 $24$~GB),
or \texttt{p4d.24xlarge} (A100 $40$~GB). Estimated wall clock per run
at the default schedule: $30$ to $60$~min on A100; $60$ to $120$~min on
A10G.

\subsection{Empirical Numbers (Measured, May 2026)}
\label{sec:empirical-numbers}

Status.  The Phase~2 PyTorch runner has executed on EC2 L40S
(g6e.xlarge).  The scaffolded predictions below have been replaced with
measurements where available.  Remaining unscaffolded items are flagged.

\subsubsection{SmolLM2-135M measurements (Exp~E2)}

\begin{table}[h]\centering\small
\caption{LoRA distillation PPL recovery on SmolLM2-135M-Instruct Q8\_0
at $r{=}8$, 1,000 WikiText-2 forward passes
(\texttt{scripts/experiment\_e2\_distill\_ppl.py}, EC2 L40S).}
\label{tab:distill-ppl}
\begin{tabular}{@{}lrrrr@{}}
\toprule
$k$ & GRC PPL & Distilled PPL & Baseline PPL & Gap Recovery \\
\midrule
256  & 98.1  & 54.3  & 6.2  & $47.7\%$ (partial) \\
512  & 7.1   & 6.0 & 6.2  & 107.1\% (over-recovery) \\
1024 & 6.5   & 6.2   & 6.2  & $98.3\%$ \\
\bottomrule
\end{tabular}
\end{table}

Key findings.
\begin{enumerate}
\item Over-recovery at $k{=}512$.  LoRA distillation recovers
      107\% of the PPL gap---the distilled model has lower PPL
      than the uncompressed baseline.  This is not ``free lunch''
      (it uses calibration data), but it demonstrates that the GRC
      residual is learnable and that the compressed representation
      can serve as a beneficial regulariser.
\item Partial recovery at $k{=}256$.  The $+1486\%$ PPL
      penalty from aggressive compression is only partially recoverable
      ($47.7\%$ gap closure).  At this rank, too much structural
      information is permanently lost.
\item Near-complete recovery at $k{=}1024$.  $98.3\%$ gap
      closure at a rank that already has low PPL penalty.
\item $\rho$ is predictive.  The recoverable-energy ratio
      $\rho{=}0.344$ (SmolLM2-135M, $k{=}256$, $r{=}8$) correctly
      predicts partial recovery; $\rho$ for $k{=}512$ is higher
      (more residual energy is within LoRA's rank-$r$ reach), explaining
      the over-recovery.
\end{enumerate}

\subsubsection{Sink-aware GRC (measured 2026-05-02)}

Sink-channel exemption protects the top-$T$ highest-magnitude columns
from GRC compression, preserving critical attention pathways while
compressing the rest.  We measured this protocol on
SmolLM2-135M-Instruct ($d{=}576$, MHA) at $k{=}512$, $T{=}32$,
comparing sink-aware GRC against vanilla GRC at the same rank.

\begin{table}[h]\centering\small
\caption{Sink-aware vs.\ vanilla GRC at $k{=}512$ on SmolLM2-135M-Instruct.}
\begin{tabular}{@{}lrrr@{}}
\toprule
Method & PPL & $\times$ baseline & Text quality \\
\midrule
Baseline (no compression) & 32.26 & $1.00{\times}$ & Paris $\checkmark$, 12×7=84 $\checkmark$ \\
Vanilla GRC, $k{=}512$  & 42.93 & $1.33{\times}$ & Paris $\checkmark$, 12×7=84 $\checkmark$ \\
Sink-aware, $T{=}32$     & 37.18 & $\mathbf{1.15{\times}}$ & Paris $\checkmark$, 12×7=88 \\
\bottomrule
\end{tabular}
\end{table}

Result: Sink-aware GRC improves PPL by 17.8\% over vanilla GRC
($42.93\to 37.18$, $\Delta{=}+5.76$ absolute), confirming the
hypothesis that protecting high-magnitude columns reduces compression
damage.  The improvement is modest but consistent: factual knowledge
is preserved (``Paris'' correct at all ranks), while arithmetic
precision is slightly degraded (12×7=88 vs.\ 84).  The 17.8\%
improvement is within the predicted range of 5--20\% for sink-channel
exemption.

\subsubsection{Large-model GRC distillation (Qwen2.5-7B)}

To test whether light distillation scales beyond 135M parameters, we
are running the identical protocol (GRC at $k{=}1024$, LoRA $r{=}8$,
500~steps on WikiText-2) on Qwen2.5-7B-Instruct ($d{=}3584$, GQA
28:4) on an EC2~L40S (24\,GB VRAM).  Result will be appended to
\path{benchmarks/llama8b_distill/} when the run completes.

\section{Threat model and limitations}
\label{sec:limits}

\subsection{What the Phase~2 run will not establish}
\begin{itemize}
\item Generalisation beyond WikiText-2. A single calibration corpus is
      not representative of all decoding workloads. The bound of
      \Cref{sec:gapbound} is corpus-conditional; the recoverable energy
      $\rho$ measured against WikiText-2 activations may differ from
      $\rho$ on a code corpus, a chat corpus, or a mathematical-reasoning
      corpus. Generalisation across corpora is out of scope of the
      Phase~2 run and is flagged for follow-up.
\item Generalisation beyond Llama-3.1-8B. The architecture
      assumptions (GQA \cite{ainslie2023gqa}, RoPE, $d{=}4096, L{=}32$)
      are baked into the trainable-parameter calculation. The
      cross-model SmolLM2 transfer check from Part~I applies to the
      projection step, not the distillation step.
\item Generalisation beyond Q4\_K\_M. Re-quantisation determinism is
      assumed; the measured PPL after Phase~3 includes a
      re-quantisation noise term that we do not characterise.
\end{itemize}

\subsection{Implementation hazards}
\begin{itemize}
\item Teacher logits at FP16 are not bit-exact across hardware. The
      student trains against the teacher run on the exact same hardware
      (a single A10G or A100 instance). Cross-instance teacher caching
      is therefore unsafe and is disabled in the runner.
\item The LoRA fit is sensitive to the AdamW $\beta_2$ default
      ($0.999$) at this corpus size; we use $\beta_2{=}0.95$ following
      QLoRA \cite{dettmers2023qlora} guidance to suppress the rare
      large-update regime.
\item Sink channels behave non-linearly under the LoRA fit if the
      fit is not warmstarted. We initialise $\Bmat_X$ to zero (so the
      adapter is identity at step $0$) and warmstart for $100$ steps
      before turning on the LoRA path, following standard LoRA
      practice.
\end{itemize}

\section{Reproducibility}
\label{sec:repro}

The Phase~1 reference is in \texttt{scripts/grc\_distill.py}. The
Phase~2 runner will live in \texttt{scripts/grc\_distill\_runner.py};
the EC2 driver in \texttt{scripts/ec2\_paperE\_distill/}. The companion
reproduction guide on the project site walks through CPU-only and GPU
paths step-by-step (see
\url{https://nagusamecs.github.io/HyperTensor/research/repro/paper-e.html}).
The pre-built W\_proj caches from Part~I are published as GitHub
Releases under the \texttt{wproj-cache-*} tag prefix and serve as the
starting point for any Phase~2 run, eliminating the calibration-free
projection step from the wall-clock budget.

The $\rho$ sweep data (Section~\ref{sec:gapbound}) is bundled in the
project's GitHub Releases under tags matching \texttt{proof-data-*}.
The 7z LZMA2 archive ($\approx$97\% size reduction) can be extracted with
\texttt{scripts/reproduce\_proof\_data.ps1}; running
\texttt{scripts/paperA\_proof/update\_paperE\_rho.py} after extraction
will regenerate the tex placeholders from the stored
\texttt{rho\_summary.json}.

A complete empirical reproduction will be published alongside the
populated results once the runner has executed. The protocol, the
predictions, and the bound in \Cref{sec:gapbound} are checked in now,
before the run, so that the run cannot retroactively become a
post-hoc rationalisation.

\printbibliography

\end{document}
